\subsection{馬其頓教會}
\subsubsection{腓立比}
一，保羅與腓立比教會的關係
\begin{itemize}%[noitemsep]
\itemsep-.25em 
\item  二次宣教
\begin{itemize}%[noitemsep]
\itemsep-.15em 
\item 保羅在特羅亞見到馬其頓異象，路加加入宣教團隊。
\item 腓立比是馬其頓頭一個城，羅馬的駐防城。他們住了幾天，安息日出城門，到了河邊禱告的地方，對聚會的婦女講道。 有一個賣紫色布匹的推雅推喇婦人呂底亞和她一家領了洗，並在家接待宣教團隊。
\item 保羅從多日喊叫的使女身上趕出汙鬼。和西拉一起被下到内監。
\item 保羅和西拉半夜禱告唱詩讚美神，地大震動，獄卒一家受洗
\item 保羅在帖撒羅尼迦時及馬其頓後，腓立比教會常供給保羅需要（腓4：15-16）
\end{itemize}
\item  四次宣教——保羅第一次坐監
\begin{itemize}%[noitemsep]
\itemsep-.25em 
\item 腓立比教會打發以巴弗提去照顧保羅（腓2：25-30）
\item 【腓立比書】由提摩太和以巴弗提送信
\end{itemize}
\end{itemize}
二，腓立比教會的問題
\begin{itemize}%[noitemsep]
\itemsep-.25em 
\item 保羅坐監是為使福音興旺，鼓勵腓立比人不要懼怕
\item 同心合意不發怨言
\item 不要靠肉體和律法稱義，而靠信心
\item 勸兩個女人尤阿迪婭和循都基要在主裡同心
\end{itemize}



\subsubsection{帖撒羅尼迦}

\subsubsection{经文}
因為底馬貪愛現今的世界，就離棄我往帖撒羅尼迦去了 (提後4：11）
保羅在帖撒羅尼迦，腓立比教會也一次兩次的打發人供給保羅的需用。（腓立比書4：16）
帖撒罗尼迦教会
                                                          

\subsubsection{I.地理和简介}
帖撒罗尼迦（简称帖）城建与主前315 年，由马其顿将军卡斯散德（Cassander）所建，並以其
妻子（亚历山大大帝的妹妹）之名「帖撒罗尼迦」命名。后来城市迅速发展，成为马其顿主要的大
都会和港口。当时有罗马帝国的军事交通大道 Agnatain way 也通过此城1.在保罗开始宣教的时候，
帖城已经有相当的规模了。当保罗第一次到帖的时候，那里已经有犹太人的会堂（徒17:1），这是
当时犹太人聚集地都会有的设施。同时在这城里也有许多希利尼人（17:4）。而帖的教会开始于保
罗的二次宣教旅程，并且保罗在第三次宣教的旅程中也拜访了这个教会。圣经中提到保罗第一次在
帖的犹太人会堂 ‘一连三个安息日，本着圣经与他们辩论’并传讲耶稣基督时，‘有些人听了劝’
其中包括许多敬虔的希利尼人和尊贵的妇女（17:2-4）。可以推论这些人就是在保罗被迫离开之后
在帖这个地方建立教会的，这样的开始离不开保罗的宣教。其中一个很重要的转折点就是发生在保
罗第二次宣教时在特罗亚看到的异象（16:6-10），一般被称之为 ‘马其顿的呼召’。因为这个异
象，保罗与同工定意从特罗亚乘船往马其顿去，而后发生了保罗在马其顿地区，腓立比，帖撒罗尼
迦，庇哩亚，和雅典，哥林多等地方的宣教事奉。我们清楚地看到是主的呼召在先，而后才有了事
工的开展。在这个过程中保罗被有嫉妒心的犹太人逼迫和追赶、驱逐，但是主一直在带领着保罗走
过每一个城市，从腓立比到哥林多都有丰富的经历神。保罗在哥林多的时候，写下了帖撒罗尼迦前
书和后书给帖的信徒和周围的人，并让他们读给所有的信徒听（帖前5）。一年零6 个月以后从坚革
哩上船往耶路撒冷去，到达耶路撒冷以后这次宣教旅程结束。在第三次宣教的行程当中，按照保罗
的计划，他走遍了马其顿一带的地方（包括帖），来到希腊，住了三个月。后因为知道犹太人要设
计害他就不走水路去耶路撒冷方向，而是走陆路回到腓立比（经过帖）从而回到特罗亚而往耶路撒
冷去（徒20）。这篇报告从分析和列举保罗事奉的榜样和帖撒罗尼迦教会的模式出发，来总结帖撒
罗尼迦教会和保罗在事奉神方面给我们留下的榜样。
\subsubsection{II.保罗在帖教会当中服侍的榜样}
保罗在建立特撒罗尼迦的过程当中树立了很好的榜样。他和他的同工西拉，提摩太一起传福音，
不是用语言，也依靠圣灵的大能并且根据对福音的确信（帖前1）。不讨好人，只取悦那位察验我们
内心的上帝。不是出于幻想、不是不良的动机、不说谄媚的话、没有藏着贪婪的念头并且确信上帝
可以作证（2:2-4）。不求称赞，没有求帖的人，也没有求别人。温柔，并且爱那里的人，愿意分享
上帝的福音，连生命也愿意给他们（2:6-9）。保罗和其同工做到了言行纯洁、公正、无可责备
（2:10），像父亲对待儿女，鼓励、安慰、劝勉帖教会的人要在生活上蒙神的悦纳，分享他的主权
和荣耀（2:11-12）。在被迫离开后，派弟兄提摩太到帖的教会，继续传扬耶稣基督。坚固和帮助他
们的信心以及预备受迫害。按照信徒的需要预先建设教会所要面对的苦难，不使魔鬼诱惑，以至于
让宣教的工作落空（3:1-5）。不光如此，保罗在写书信的时候也提到，他感谢上帝，常常祷告，也
祝福当地教会（帖前1:1-2；5:23-28;帖后1:1-2;3:16-17）并且因为帖的教会常常喜乐，受安慰和
鼓舞（帖前3:7-13）。保罗提到的这些事迹是他在帖城时候教会的弟兄姐妹亲眼目睹见证的，帖的
弟兄姐妹在传阅和诵读这两封信的时候一定会受到鼓舞和激励。
\subsubsection{III.帖撒罗尼迦教会的榜样}
而帖的教会也成为全地的榜样，使保罗因着他们的信心，大得安慰和鼓舞。帖的教会的信徒在
生活上讨神的喜欢（帖前4:1）。他们把所信的实行出来，以爱心辛劳工作，坚守对我们的主耶稣基
督的盼望（1:3）。效法使徒，效法主，虽遭受大难，仍然在圣灵所赐的喜乐中领受信息，成为马其顿和亚该亚所有信徒的模范（1:6-7）。主的信息从这里传开，对上帝的信心也传到远近各处。他们
接待使徒（宣教士），离弃偶像，归向上帝，事奉神，并且盼望耶稣再来（1:8-10）。他们领受那
确确实实出自上帝的信息，而上帝也在信的人当中工作（2:13）。有信心和爱心，也想念使徒，迫
切希望见到使徒（和使徒想念他们一样，有同样的心思和意念）（3:6）。因此信心格外增加；彼此
相爱的心充足；受患难存忍耐和信心，而这是神公义审判的明证、算配得神的国，并且为这国受苦
（帖后1:3-5）。后来在保罗的宣教团队当中出现了帖人亚里达古和西公都，以及马其顿其他地区的
同工，庇哩亚的毕罗斯的儿子所巴特，特庇人该犹，从而看到了在马其顿地区宣教和建立教会所结
的果子（徒20:4）。

\subsubsection{IV.保罗和帖的教会的挑战}
但是使徒保罗和帖的教会在这个过程中也遇到很大的挑战。他们最大的挑战始终是犹太人的迫
害，包括保罗被阻挠不可向外邦人传福音，这跟犹太地区上帝的各教会，就是属耶稣基督的信徒们
所遭遇的一样。他们受同胞（犹太人）的迫害，犹太人杀害先知，驱逐使徒，冒犯上帝，跟全人类
为敌（帖前2:14-18）。
其次也有世界的挑战，帖是繁华的大城，其中充斥着世界的影响，很可能会贪爱世界（提后
4:10）。根据帖的教会的实际情况，保罗在书信当中对他们重新申明真道，可以看作是对于他们的
劝勉、鼓励和教导。同时也看出帖的教会的信徒所面对的挑战以及因此产生的疑惑。保罗教导说要
圣洁，要绝对没有淫乱的事。依照圣洁合宜的方法控制自己的身体，而不是凭借情欲，像不认识上
帝的异教徒一样。不可以对不起同道或占他便宜的行为，因为这样的行为主必惩罚。那些拒绝过圣
洁生活的，不是拒绝人，而是拒绝神。要彼此相爱，要更加努力，立志安分守己，亲手做工来维持
自己的生活，从而得到非信徒的尊重，同时也不用依赖别人的供应。论到已经死了的人，应当有盼
望而不忧伤。要相信死而复活，要一同跟主相会，跟主在一起（帖前4:2-17）。主再来的日子不需
要人家告诉你们在什么时候，那是我们不知道的，只要时刻保持警醒（贴前5:1-10）。要彼此鼓励、
彼此帮助、尊重按主所选召劝诫你们的人，有敬意和爱心对待这些人，自己也要和睦（5:11-13）。
要警告懒惰的人，鼓励灰心的人，扶助软弱的人，耐心待每一个人，要谨慎，不可以以恶报恶，彼
此关心，为别人的好处着想。要常常喜乐，常常祷告，凡事感谢。不可以抑制圣灵的工作，不可轻
视信息的传讲。详细察验每一件事，持守美善的，弃绝邪恶的（5:14-22）。而神是公义的，那些受
患难的人必得平安（帖后1:3-12）。你们当中有假冒的使徒和教师，不要被诱惑，因为必有离道反
教的事，就是沉沦之子显露出来。但是要在真道上站立得住（2 章）。劝勉他们要爱神，并且学基督
的忍耐（帖后3:5）。要安静做工，行善不可丧志，效法使徒，按规矩而行（3:6-15）。
所以通过保罗对于这些真理的阐述，可以推想帖的教会内部当时可能面对包括如何过圣洁生活
的挑战；对于死前是否能等到主再来的疑惑；以及假教师或者先知的教导；并且教会里面也有灰心
丧气的人，不愿做工。但是帖的教会的信徒们显然仍然是有渴慕真道，彼此相爱，存爱心，信心以
及盼望过生活。为他们能够对抗来自外部的逼迫和诱惑，以及来自的内部的搅扰，仍能够追求真道
奠定了坚实的基础，以致成为全地的榜样。
\subsubsection{V.保罗和帖撒罗尼迦教会的关系}
保罗在帖的教会停留的时间可能总共都不长，记载了他第一次停留了三个安息日，之后被赶了
出来，后面就没有记载，但是他去的次数很多。第二次宣教他去过一次，第三次宣教他很可能去过
两次。通过对于保罗和帖的教会交往的观察，他因为帖的教会渴慕真道的有点并且成为全地的结果
而大受安慰，所以他多次愿意去到帖的教会并且在书信中表达他的喜乐和受安慰之情。但是就像保
罗自己也时常遇到挑战一样，他也提醒了许多帖的信徒要注意的事情。保罗会担心帖的教会收到诱
惑和假的教导而不能持守真道，以至于在真道上站立不住，所以常常要为他们祷告，记念他们和想
见他们。他时常鼓励帖的信徒要更加竭力地过圣洁的生活。并且在自己不能去的时候，因为担心他
们的状态而派同工提摩太去帖撒罗尼迦的教会，要多了解他们的消息。可以说整个马其顿的宣教工
作是从马其顿的异象开始，保罗和同工知道这地区的人需要帮助，并且定意要去。他们在宣教的过
程当中，在腓立比的时候已经发现了主的预备，那里有个祷告的地方（徒16:13），并且有卖紫色布
的妇人吕底亚和狱卒一家信主，坚固了宣教团队的信心。在帖撒罗尼迦有犹太人的会堂，保罗进去
辩论，并且有妇人和希利尼人信主，但是保罗受到犹太人的迫害。犹太人所嫉妒和弃绝的，外邦人
却接受了，保罗看到主的道在全地传开应该感到欣慰。从开始帖撒罗尼迦教会是保罗播种传福音的
地方，然后这个教会成为了全地的榜样，同心追求真道，最后同保罗一同从这个地区走出去往亚细
亚的同工有7 位之多，而有两位是来自于帖撒罗尼迦。保罗不光在这里播洒下了种子，并且也欢呼
收割。在彼此面对挑战的时候保罗和帖的教会也成为了很好的互相祷告的对象，不光是保罗常常为
帖的教会祷告，他也会把自己和同工团的祷告事项分享给帖的人（帖后3 章）。这样我们看到的不
是教导与被教导的关系，而是弟兄姐妹之间彼此亲密无间，互相分担，一同求告神的关系。以及保
罗在帖的时候，辛苦劳作，并且用温柔的心和劝勉的话，有母亲和父亲的心肠（帖前2:7-12）。

一，保羅與帖撒羅尼迦教會的關係
\begin{itemize}%[noitemsep]
\itemsep-.25em 
\item 保羅在帖撒羅尼迦猶太會堂一連三個安息日，本著聖經講解耶穌
\item 有些人附從保羅和西拉，有許多虔敬的希臘人和不少尊貴的婦女
\item 不信的猶太人心裡嫉妒，招聚市井匪類把耶孫和幾個弟兄拉到地方官，官長取了耶孫和其餘之人的保狀
\end{itemize}
\newpage
二，帖前要解決的問題
\begin{itemize}%[noitemsep]
\itemsep-.25em 
\item 可能對保羅快速的離開他們有誤解“以爲保羅怕死”？
\item 成為聖潔，遠避淫行，不要沾染汙穢
\item 因不明白復活而為睡了的人而憂傷
\item 勸他們驚醒，等候主的再來
\end{itemize}
三，帖后要解決的問題
\begin{itemize}%[noitemsep]
\itemsep-.25em 
\item 在患難和逼迫中要站立得穩
\item 冒充保羅的名口傳或者寫信說耶穌的日子已經到了
\item 有人不做工，反倒專管閑事
%\item  4 勸他們驚醒，等候主的再來
\end{itemize}

\subsubsection{庇哩亞——徒17：10-14}

\begin{itemize}%[noitemsep]
\itemsep-.25em 
\item 帖撒羅尼迦的地方官取了耶孫等人的保狀，就釋放了他們，眾人隨即在夜間打發保羅和西拉往庇哩亞去。
\item 保羅和西拉到了，就進入庇哩亞猶太人的會堂。
\item 這地方的人賢於帖撒羅尼迦的人，甘心領受這道，天天考查聖經，要曉得這道是與不是。 
\item 所以他們中間多有相信的，又有希利尼尊貴的婦女，男子也不少。
\item 帖撒羅尼迦的猶太人知道保羅又在庇哩亞傳神的道，也就往那裡去，聳動攪擾眾人。
\item 弟兄們便打發保羅往海邊去，西拉和提摩太仍住在庇哩亞。 
\end{itemize}

\subsection{亞該亞教會}
\subsubsection{林前，林後}
一，保羅與哥林多教會的關係
\begin{itemize}%[noitemsep]
\item 保羅第二次宣教
\begin{itemize}%[noitemsep]
\itemsep-.25em 
\item 保羅在哥林多遇見猶太同行，製造帳篷為業的夫婦亞居拉和妻百基拉從羅馬來。保羅就投奔了他們
\item 每逢安息日保羅在會堂勸化猶太人和希臘人，遭猶太人抗拒
\item 保羅到了靠近會堂提多.猶士都的家， 管會堂的基利司布和全家都信了主，許多哥林多人聽了就相信、受洗。司提反一家是亞該亞初結的果子，並且他們專以服侍聖徒為念（林前16：15）
\item 保羅在那裡住了一年零六個月，將神的道教訓他們
\item 猶太人在亞該亞（省）的方伯（省長）迦流面前控告保羅,眾人毆打管會堂的所提尼
\item 保羅又住了多日和百基拉、亞居拉坐船到以弗所
\item 保羅離開以弗所後，亞波羅到以弗所遇到百基拉、亞居拉。亞波羅又去哥林多繼續作工，他在眾人面前，極有能力，駁倒猶太人，引證聖經，證明耶穌是基督（徒18：27，28，19：1）
\end{itemize}
\item 保羅第三次宣教
\begin{itemize}%[noitemsep]
\itemsep-.25em 
\item 保羅曾給哥林多人寫了第一封信，遺失了（林前5：9-11）
\item 革來氏家中有人去以弗所，報告哥林多教會紛爭的事情（林前1：11）
\item 保羅打發提摩太由以弗所起程，去馬其頓和哥林多（林前4：17，16：10），但提摩太可能只到了馬其頓（徒19：22）
\item 哥林多人寫信給在以弗所的保羅（林前7：1）。可能是托司提反、福徒拿都和亞該古帶信給保羅（林前16：17）
\item 保羅和所提尼（林前1：1）在以弗所（林前16：8）於五旬節前寫了林前(第二封信)，給哥林多，林前第7章以後是回答哥林多信中的問題。亞波羅此時與保羅同在，保羅曾勸他去哥林多，但遭拒絕（林前16：12）。
\item 保羅可能第二次去了哥林多（林后12：14保羅自己提到要第三次造訪哥林多）
\item 哥林多教會發生新困難，有人對保羅作使徒的職權起了懷疑，說他的信（可能指林前），又沉重、又厲害，是氣貌不揚，言語粗俗，一個詭詐的人（林後10：7，10，11：5，6，12：16，17）。
\item 保羅寫給哥林多教會第三封「流淚的信」（林後2：3、4），該信已失落。打發提多到哥林多（林后7：15，8：16-18，12：18）並囑咐他由馬其頓到特羅亞在那裡回見保羅（林後2：12，13）。
\item 保羅在馬其頓（林後7：5，6）得到提多的報告，大得安慰鼓勵，他和提摩太（林后1：1）在馬其頓作哥林多後書，打發二位弟兄一同送往哥林多（林後8：1，16-24）
\item 保羅第三次到哥林多，在那裡住了三個月（徒20：2，3），由此寫羅馬書
\end{itemize}
%\item  4 
\end{itemize}

二，林前要解決的問題
\begin{itemize}%[noitemsep]
\itemsep-.25em 
\item 解決保羅從信徒口中聽到的哥林多的問題（1-6章）
\begin{itemize}
\itemsep-.25em 
\item 分黨-中間有紛爭。
\item 淫亂-有人收了他的繼母。
\item 告狀-彼此相爭，告在不信主的人面前
\end{itemize}
\item  解決哥林多教會信中的問題（7-16章）
\begin{itemize}
\itemsep-.25em 
\item 婚姻問題
\item 身份
\item 吃祭偶像之物
\item 女人蒙頭
\item 混亂聖餐
\item 屬靈恩賜
\item 死人復活
\item 給耶路撒冷教會捐錢的事
\end{itemize}
\end{itemize}
\newpage
三，林後要解決的問題
\begin{itemize}%[noitemsep]
\itemsep-.25em 
\item 保羅為自己的辯解（1-6章）
\begin{itemize}
\itemsep-.25em 
\item 兩次修改行程的原因
\item 自己成為新約的執事
\item 在各樣事上表明自己是神的用人
\end{itemize}
\item 勸哥林多樂意捐款給耶路撒冷教會（7-9章）
\item 警告哥林多自己要第三次再去哥林多（10-13章）
\begin{itemize}
\itemsep-.25em 
\item 保羅與假使徒的區別
\item 第三次要帶著使徒權柄再去哥林多
\end{itemize}
\end{itemize}
%\subsubsection{林后}
\subsubsection{堅革哩}
\begin{itemize}%[noitemsep]
\itemsep-.25em 
\item 保羅因為許過願，就在堅革哩剪了頭髮，然後到以弗所
\item 保羅第三次到哥林多时，托堅革哩教會的女執事菲比帶羅馬書給羅馬教會
\end{itemize}

\subsection{羅馬書}
\subsubsection{羅馬教會介紹}
\begin{itemize}%[noitemsep]
\itemsep-.25em 
\item 第三次宣教
\begin{itemize}%[noitemsep]
\itemsep-.15em
\item 保羅写【羅馬書】时自己沒有到過羅馬教會（羅15：23）
\item 保罗盼望罗马教会可以蒙罗马教会送行，去西班牙传福音（羅15：22-28）
\item 保罗在哥林多写【羅馬書】后，托堅革哩教會中的女執事非比带信（羅16：1-2）
\item 【羅馬書】由德丟代筆，同工:提摩太,路求,耶孫,所西巴德,該猶,管銀庫的以拉都,括土（羅16：21-24）
\item 保罗写【羅馬書】时，百基拉和亞居拉住在羅馬，並且在他們家中有教會（羅16：3-4）
\item 羅馬教會的成員由猶太人和外邦人組成，有弟兄也有姐妹（羅16：5-15）
\end{itemize}
\item 第一次在罗马做监——四次宣教
\begin{itemize}%[noitemsep]
\itemsep-.25em
\item 保罗在耶路撒冷被抓后，由罗马人押着坐船到罗马（徒27-28）
\item 在罗马自己租的屋子住了两年，天天向人（包括犹太人）传讲耶稣（徒28：30）
\item 虽然在罗马监狱的困锁中，保罗将福音传遍了御营军（腓1：13）
\end{itemize}
\item 第二次在罗马做监
\begin{itemize}%[noitemsep]
\itemsep-.25em
\item 保罗后来第二次被囚在罗马监狱，并死于罗马（提后4：6）
\end{itemize}
\end{itemize}

\subsubsection{羅馬書的簡綱}
\begin{itemize}%[noitemsep]
\itemsep-.25em 
\item 神是審判世界的主 1
\begin{itemize}%[noitemsep]
\itemsep-.25em 
\item 保羅蒙召做使徒的職分 1：1-17
\item 世人都犯了罪 1：18-32
\end{itemize}
\item 保羅寫給猶太人的話 2-8
\begin{itemize}%[noitemsep]
\itemsep-.25em 
\item 猶太人自認為有律法和割禮就論斷人 2
\item 神賜給猶太人律法的目的 3
\item 信心和律法的討論 4-5
\item 靠聖靈在基督裡得新生命 6-8
\end{itemize}
\item 保羅寫給外邦人的話 9-15 
\begin{itemize}%[noitemsep]
\itemsep-.25em 
\item 真以色列人的定義 9
\item 因著以色列人的不信，外邦人得著救恩，最後以色列全家都要得救  10-11
\item 在基督裡的新生命 12-15
\end{itemize}
\item 最後的問安 16
\end{itemize}

\subsubsection{羅馬書要解決的問題}
\begin{itemize}%[noitemsep]
\itemsep-.25em 
\item 猶太人的問題
\begin{itemize}%[noitemsep]
\itemsep-.25em 
\item 猶太人自認為有律法和割禮就論斷人 

\begin{itemize}%[noitemsep]
\itemsep-.25em 
%\item 自恃熟悉律法而论断人 2:1-11
\item 对有无律法的审判标准不清楚 2:12-16
\item 自恃熟悉律法却知法犯法 2:17-24
\item 未受割礼却全守律法的人要审判受过割礼却犯律法的人 2:25-27
\item 真割礼和真犹太人的定义 2:28-29
\end{itemize}

\item 因律法稱義還是靠信心 
\begin{itemize}%[noitemsep]
\itemsep-.25em
\item 犹太人的不信不能废掉神的信 3:1-4
\item 犹太人的不信反而显出神的荣耀，神发怒不合理（犹太人虽然不信，但他们守住了律法，所以神不应当因为他们不信而发怒） 3:5-8
\item 对犹太人有罪的概念不清楚 3:9-20 
\item 神在律法以外显明的义 3:9-20 
\item 因信耶稣而坚固律法 3:21-31
\end{itemize}
\end{itemize}


\item 外邦人的問題
\begin{itemize}%[noitemsep]
\itemsep-.25em 
\item 以色列人的不信說明神的話落空了嗎？
\item 神棄絕以色列人了嗎？ 
\end{itemize}
\end{itemize}
